\section{Discussion and Limitations}

The results support three conclusions. First, the folding-layout control shows that the sample ordering matters beyond the increase in input height: phase-interleaved folding beats a contiguous reshape in 25 of 27 selection-metric comparisons. The geometry rule is therefore a reproducible prior, not a claim that one fold factor is optimal for every task. Second, visual initialization contributes substantially under the matched downstream budget: the selection metric improves over random initialization in 34 of 36 backbone--dataset comparisons (balanced accuracy improves in 33). Third, the effect transfers across CNN and Transformer backbones, although the strongest backbone depends on the task.

The benchmark comparison gives the same qualified picture. ConvNeXt-Tiny exceeds the strongest included published EEG foundation reference on both main metrics in 8 of 12 tasks. Under the matched downstream reruns, at least one visual backbone exceeds both CBraMod and REVE-Base on balanced accuracy and the selection metric in 8 of 12 datasets. The exceptions are SHU-MI, Mumtaz2016, PhysioNet-MI, and SEED-V, where REVE-Base remains stronger on at least one metric. These reruns use five seeds for every model--dataset pair, but they do not reproduce every original training stage: CBraMod uses a uniform learning rate, REVE omits its published warm start and additional training procedures, and a few datasets require documented fallback input, readout, or learning-rate settings. They should therefore be read as protocol-controlled comparisons; the published values remain the external benchmark references.

Taken together, the findings establish visual pretraining with geometry alignment as a strong and reproducible EEG baseline that does not require an EEG-specific pretraining stage. This is useful when EEG pretraining data or compute are limited and provides a necessary control for attributing downstream gains to EEG-specific pretraining. The result is conditional on the evaluated full fine-tuning setting: EEG-specific models remain stronger on several tasks, and these experiments do not establish advantages for few-shot, frozen, or cross-montage transfer.

Several limitations remain. Most task-specific EEG baselines and the main-table foundation-model numbers are published references rather than paired local reruns, despite the additional matched CBraMod and REVE reruns in the appendix. The folding adapter preserves stored channel order but does not use physical electrode coordinates; robustness to channel permutation, missing electrodes, and cross-montage transfer is untested. Finally, fold factor selection is fixed by a geometry rule, and the study does not report paired significance tests or pretraining-efficiency comparisons. These limitations define directions for future work.
